\section{Second Collapse: Feedback Stops Testing the Claim}

An image arrives before the user has learned to draw. A conversation offers attention without requiring another person's time. A game supplies a contest with carefully managed stakes. These experiences can be enjoyable, useful, and educational. Their effects depend on which feedback they preserve.

\subsection{A Feeling and the Process That Once Produced It}

Technical systems can separate a reward from some of the effort, uncertainty, or obligation that previously accompanied it. A finished artifact can arrive without practice in making it. An agreeable interaction can arrive without the negotiation of a reciprocal relationship.

The feeling is real. The question is what else develops with it.

A drawing tool can teach composition when the user compares, revises, and learns. It can also deliver an image while leaving those skills unchanged. A conversational system can support a person through a difficult moment. It can also occupy time that would otherwise sustain mutual obligations. The interface alone does not determine the result.

Fiction and play have always offered experience at a distance from its full consequences. Their value includes precisely that distance: one can explore a dangerous choice without living its worst outcome. Training simulators turn the same separation into a practical instrument.

The crucial connection is transfer. Which judgment is being trained, by which feedback, and against what independent standard? A simulator can alter a person's skills even though its aircraft is virtual. A system can produce intense feedback while teaching little about the world outside itself.

The second collapse occurs when the feedback ceases to test the claim for which it is being used. Enjoyment is used as evidence of competence. Reassurance is used as evidence of a relationship's reciprocity. A convincing explanation is used as evidence of a mechanism.

\subsection{Can Reality Interrupt the Representation?}

Baudrillard's account of simulacra follows changing relations between signs and what they purport to represent \cite{baudrillard1981simulacres}. Borges's map as large as the territory provides a related image of representation overwhelming its referent \cite{borges1946rigor}.

For this book, the operational test is straightforward: what can interrupt the representation?

An experiment can contradict an explanation. A load test can defeat a design. A source record can disprove a citation. A user can reject a proposed interpretation of their life. These are points at which the generating system does not control the answer.

Surface quality and external support can vary independently. A rough field note can preserve excellent evidence. A polished generated report can have none. There is no general law that signs become more convincing as their reality connection weakens; that combination is a dangerous possibility, not a mandatory sequence.

Goodhart and Campbell identify related problems when feedback becomes a target \cite{goodhart1975problems,campbell1979assessing}. A department learns to improve the reported measure. A researcher learns the benchmark. A platform learns which response receives a favorable rating. Each can increase the score while weakening its connection to the intended outcome.

The same issue appears inside machine learning. Training on evaluation items can inflate apparent generalization \cite{sainz2023nlp}. Recursively reusing model outputs can reduce information about less common patterns under particular conditions \cite{shumailov2024ai}. Carefully selected synthetic data can also be useful. The causal variables include quality, diversity, coverage, filtering, and access to fresh evidence.

Provenance tools supply one concrete response. C2PA's version 2.4 specification provides mechanisms for associating signed claims and content history with an asset \cite{c2pa2026spec}. Its validation concerns the provenance assertions and their integrity. It does not certify the truth of every depicted event or the honesty of a signer.

That distinction is an example of a healthy interface. A technical standard says which claim its check supports. Users then know which additional checks remain necessary.

The problem becomes severe when a representation environment insulates itself from contrary evidence. A ranking supplies the examples used to judge the ranking. A model supplies the explanation used to excuse its errors. An organization evaluates its performance using reports generated entirely inside the same reward structure.

An independent interruption point keeps such a system learnable. It can be a physical measurement, a fresh test, an external source, an appeal, or a person with authority to say that the representation has missed what matters.

\subsection{Who Arranges the Things We Want?}

Feedback changes desires as well as beliefs. Girard's triangular account places a mediator between subject and object: people often learn what to want by observing what others want \cite{girard1961mensonge}. Rivalry can intensify when model and observer compete in the same field \cite{girard1972violence}.

Digital platforms change the available mediators. A feed selects which lives, achievements, bodies, purchases, and opinions appear nearby. It can omit the preparation, financing, failure, and selection that produced the display. Viewers encounter a result without its cost structure.

Social-comparison research provides a narrower empirical framework \cite{festinger1954theory}. Studies of passive browsing and well-being investigate effects that differ by behavior, person, and setting \cite{verduyn2015passive}. More mediators do not automatically make everyone want the same thing. They can diversify aspirations, intensify local competition, or do both.

The pertinent design question is how selection changes the distribution of apparent possibilities. Which alternatives become visible? Which are rewarded? Which costs disappear from view? A person's choice occurs inside that presentation.

The temptation to treat material life as a dispensable inconvenience then has a long intellectual echo. Some ancient Gnostic traditions described a damaged material world and a route of escape through special knowledge. The analogy here concerns a preference for representation over the burden of living its counterpart. The historical traditions were diverse; a modern screen habit is not their direct continuation.

The stronger claim is practical. If a service supplies an attractive experience more cheaply than the activity it resembles, people have reason to shift time toward the service. Whether that shift enriches or narrows a life depends on learning, relationships, health, opportunity, and the alternatives actually available.

The test of authenticity is therefore not how painful an experience was, how many materials it consumed, or whether a computer participated. It is whether the claim made for the experience remains connected to something capable of correcting it.

When an experience is sold as competence, test transfer. When it is sold as evidence, inspect the source. When it is sold as a relationship, inspect the obligations and control. The surface can remain delightful while those questions stay open.
